%%--------------------------------------------------------------------
%1--------------------------------------------------------------------
\vspace*{\fill}
\begin{glyph}{Achieve (達)}{achieve}\label{glyph:achieve}
Achieve; Attain.\\\hline
When used as a suffix, this kanji also serves to indicate the plural, although the second compound below can refer to only one person as well.
\end{glyph}

\begin{comps}
\comprtwocone & 土 + 羊 + 辶 \\\hline
\comprthreecone & Earth + Sheep + Road/Walk \\\hline
\end{comps}

\begin{remglyph}
\hooglig{Achieve} \remcom{dirty} \remcom{sheep} by letting them \remcom{walk}.\\\hline
\end{remglyph}

\begin{prons}
\pronsrtwocone & \textbf{タツ (TATSU)} & \textbf{たち (tachi)}\\\hline
\rye{3}{The second example shows how the kun-yomi can become voiced; this reading never appears in the first position.}
\pronsrthreecone & --- & --- \\\hline
\end{prons}

\begin{remprons}
\hooglig{Achieve} \comon{tuts} of disapproval due to your \comkun{touchiness}.\\\hline
\end{remprons}

%{\middel}
% &  \mbox{()}\\\hline
\begin{somme}
上達 & upper + achieve = \textit{improvement; progress} & ジョウ{\middel}タツ \mbox{(J\={O}{\middel}TATSU)}\\\hline
友達 & friend + achieve = \textit{friend(s)} & とも{\middel}だち \mbox{(tomo{\middel}dachi)}\\\hline
\notyet{発達} & discharge + achieve = \textit{correspondence} & ハッ{\middel}タツ \mbox{(HAT{\middel}TATSU)}\\\hline
\end{somme}
\end{multicols}
\vhrulefill{5pt}
%%--------------------------------------------------------------------
%2--------------------------------------------------------------------
\begin{multicols}{2}
\begin{glyph}{Continue (続)}{continue}\label{glyph:continue}
Continue.
\end{glyph}

\begin{comps}
\comprtwocone & 糸 + 売 \\\hline
\comprthreecone & Thread + Sell (\ref{glyph:sell}) \\\hline
\end{comps}

\begin{remglyph}
\hooglig{Continue} \remcom{selling} \remcom{thread}.\\\hline
\end{remglyph}

\begin{prons}
\pronsrtwocone & \textbf{ゾク (ZOKU)} & \textbf{つづ (tsuzu)}\\\hline
\pronsrthreecone & --- & --- \\\hline
\end{prons}

\begin{remprons}
\hooglig{Continuation} of the \comkun{stew zucchini} \comon{tribe}.\\\hline
\end{remprons}

%{\middel}
% &  \mbox{()}\\\hline
\begin{somme}
続く (intr.) & \textit{continue} & つず{\middel}く \mbox{(tsuzu{\middel}ku)}\\\hline
続ける (tr.) & \textit{continue} & つず{\middel}ける \mbox{(tsuzu{\middel}keru)}\\\hline
手続き & hand + continue = \textit{formalities; procedure} & て　つず{\middel}き \mbox{(te tsuzu{\middel}ki)}\\\hline
\notyet{接続} & join + continue = \textit{connection; joining} & セツ{\middel}ゾク \mbox{(SETSU{\middel}ZOKU)}\\\hline
\end{somme}
\end{multicols}
\vspace*{\fill}
\pagebreak
%%--------------------------------------------------------------------
%3--------------------------------------------------------------------
\vspace*{\fill}
\begin{multicols}{2}
\begin{glyph}{Weak (弱)}{weak}\label{glyph:weak}
Weak.
\end{glyph}

\begin{comps}
\comprtwocone & 弓 + 冫 \\\hline
\comprthreecone & Bow + Ice \\\hline
\end{comps}

\begin{remglyph}
\hooglig{Weak} \remcom{ice} \remcom{bows}.\\\hline
\end{remglyph}

\begin{prons}
\pronsrtwocone & \textbf{ジャク (JAKU)} & \textbf{よわ (yowa)}\\\hline
\pronsrthreecone & --- & --- \\\hline
\end{prons}

\begin{remprons}
\hooglig{Weakling} \comon{jucked} \comkun{your watch} away.\\\hline
\end{remprons}

%{\middel}
% &  \mbox{()}\\\hline
\begin{somme}
弱い & \textit{weak} & よわ{\middel}い \mbox{(yowa{\middel}i)}\\\hline
\notyet{弱点} & weak + point = \textit{weak point; shortcoming} & ジャク{\middel}テン \mbox{(JAKU{\middel}TEN)}\\\hline
\notyet{弱肉強食} & weak + meat + strong + eat = \textit{survival of the fittest} & ジャク{\middel}ニク{\middel}キョウ{\middel}ショク (JAKU{\middel}NIKU{\middel}K\={O}{\middel}SHOKU)\\\hline
\end{somme}
\end{multicols}
\vhrulefill{5pt}
%%--------------------------------------------------------------------
%4--------------------------------------------------------------------
\begin{multicols}{2}
\begin{glyph}{Short (短)}{short}\label{glyph:short}
Short.
\end{glyph}

\begin{comps}
\comprtwocone & 矢 + 豆 \\\hline
\comprthreecone & Arrow + Bean/A Kind of a Vessel \\\hline
\end{comps}

\begin{remglyph}
\hooglig{Short} \remcom{arrows} in a \remcom{kind of a vessel}.\\\hline
\end{remglyph}

\begin{prons}
\pronsrtwocone & \textbf{タン (TAN)} & \textbf{みじか (mijika)}\\\hline
\pronsrthreecone & --- & --- \\\hline
\end{prons}

\begin{remprons}
The \hooglig{short} \comkun{midget cart} was \comon{tons} of fun.\\\hline
\end{remprons}

%{\middel}
% &  \mbox{()}\\\hline
\begin{somme}
短い & \textit{short} & みじか{\middel}い \mbox{(mijika{\middel}i)}\\\hline
最短 & most + short = \textit{the shortest} & サイ{\middel}タン \mbox{(SAI{\middel}TAN)}\\\hline
\notyet{短期} & short + period = \textit{short-term} & タン{\middel}キ \mbox{(TAN{\middel}KI)}\\\hline
\end{somme}
\end{multicols}
\vspace*{\fill}
\pagebreak
%%--------------------------------------------------------------------
%5--------------------------------------------------------------------
\vspace*{\fill}
\begin{multicols}{2}
\begin{glyph}{Resign (辞)}{resign}\label{glyph:resign}
Resign, Quit.\\\hline
A second meaning relates to phrases and words; this sense is reflected in the final two compounds.
\end{glyph}

\begin{comps}
\comprtwocone & 舌 + 辛 \\\hline
\comprthreecone & Tongue + Spicy \\\hline
\end{comps}

\begin{remglyph}
\hooglig{Resign} your \remcom{tongue} to the \remcom{spiciness}.\\\hline
\end{remglyph}

\begin{prons}
\pronsrtwocone & \textbf{ジ (JI)} & \textbf{や (ya)}\\\hline
\pronsrthreecone & --- & --- \\\hline
\end{prons}

\begin{remprons}
\hooglig{Quit} \comkun{young} \comon{jigs}.\\\hline
\end{remprons}

%{\middel}
% &  \mbox{()}\\\hline
\begin{somme}
辞める & \textit{resign; quit} & や{\middel}める \mbox{(ya{\middel}meru)}\\\hline
辞書 & resign (phrase) + write = \textit{dictionary} & ジ{\middel}ショ \mbox{(JI{\middel}SHO)}\\\hline
お世辞 & society + resign (phrase) = \textit{flattery; compliment} & お　セ{\middel}ジ \mbox{(o SE{\middel}JI)}\\\hline
\end{somme}
\end{multicols}
\vhrulefill{5pt}
%%--------------------------------------------------------------------
%6--------------------------------------------------------------------
\begin{multicols}{2}
\begin{glyph}{Exemption (免)}{exemption}\label{glyph:exemption}
Exemption.  Permission.
\end{glyph}

\begin{comps}
\comprtwocone & 免 \\\hline
\comprthreecone & Exemption/Dismissal \\\hline
\end{comps}

\begin{remglyph}
Uncommon Kanji radical.\\\hline
\end{remglyph}

\begin{prons}
\pronsrtwocone & \textbf{メン (MEN)} & --- \\\hline
\pronsrthreecone & --- & まぬか (manuka) \\\hline
\end{prons}

\begin{remprons}
\hooglig{Exempt} \comon{men} from \ucomkun{manure customs}.\\\hline
\end{remprons}

%{\middel}
% &  \mbox{()}\\\hline
\begin{somme}
免じる & \textit{to exempt; dismiss} & メン{\middel}じる \mbox{(MEN{\middel}jiru)}\\\hline
\end{somme}
\end{multicols}
\vspace*{\fill}
\pagebreak
%%--------------------------------------------------------------------
%7--------------------------------------------------------------------
\vspace*{\fill}
\begin{multicols}{2}
\begin{glyph}{Song (歌)}{song}\label{glyph:song}
Song, singing.\\\hline
In a few compounds, this character can also mean poetry.
\end{glyph}

\begin{comps}
\comprtwocone & 可 + 可 + 欠 \\\hline
\comprthreecone & Possible (\ref{glyph:possible}) + Possible (\ref{glyph:possible}) + Lack \\\hline
\end{comps}

\begin{remglyph}
\hooglig{Songs} are \remcom{possibly} \remcom{lacking}.\\\hline
\end{remglyph}

\begin{prons}
\pronsrtwocone & \textbf{カ (KA)} & \textbf{うた (uta)}\\\hline
\rye{3}{The on-yomi is the more common reading.}
\pronsrthreecone & --- & --- \\\hline
\end{prons}

\begin{remprons}
A \hooglig{song} that \comon{cusses} \comkun{Utah}.\\\hline
\end{remprons}

%{\middel}
% &  \mbox{()}\\\hline
\begin{somme}
歌 & \textit{song; poem} & うた \mbox{(uta)}\\\hline
歌手 & song + hand = \textit{singer} & カ{\middel}シュ \mbox{(KA{\middel}SHU)}\\\hline
国歌 & country + song = \textit{national anthem} & コッ{\middel}カ \mbox{(KOK{\middel}KA)}\\\hline
\end{somme}
\end{multicols}
\vhrulefill{5pt}
%%--------------------------------------------------------------------
%8--------------------------------------------------------------------
\begin{multicols}{2}
\begin{glyph}{Send (送)}{send}\label{glyph:send}
Send.
\end{glyph}

\begin{comps}
\comprtwocone & 艹 + 大 + 辶 \\\hline
\comprthreecone & Grass/Herb/Plant + Big/Large + Road/Walking\\\hline
\end{comps}

\begin{remglyph}
\hooglig{Send} \remcom{big} \remcom{herbs} to the \remcom{road}.\\\hline
\end{remglyph}

\begin{prons}
\pronsrtwocone & \textbf{ソウ (S\={O})} & \textbf{おく (oku)}\\\hline
\rye{3}{Although there are a few compound verbs making use of the kun-yomi, the on-yomi is the more common reading.}
\pronsrthreecone & --- & --- \\\hline
\end{prons}

\begin{remprons}
\hooglig{Send} the \comon{sore} \comkun{occupants} away.\\\hline
\end{remprons}

%{\middel}
% &  \mbox{()}\\\hline
\begin{somme}
送る & \textit{to send} & おく{\middel}る \mbox{(oku{\middel}ru)}\\\hline
見送る & see + send = \textit{see (someone) off} & み{\middel}おく{\middel}る \mbox{(mi{\middel}oku{\middel}ru)}\\\hline
\notyet{送別} & send + different = \textit{farewell; send off} & ソウ{\middel}ベツ \mbox{(S\={O}{\middel}⎄BETSU)}\\\hline
\end{somme}
\end{multicols}
\vspace*{\fill}
\pagebreak
%%--------------------------------------------------------------------
%9--------------------------------------------------------------------
\vspace*{\fill}
\begin{multicols}{2}
\begin{glyph}{Pay (払)}{pay}\label{glyph:pay}
Pay.\\\hline
A secondary meaning relates to the idea of `driving/sweeping away' something.
\end{glyph}

\begin{comps}
\comprtwocone & 手扌才 + 厶 \\\hline
\comprthreecone & Hand/Actions + Self/Private \\\hline
\end{comps}

\begin{remglyph}
\hooglig{Pay} for your \remcom{actions} by \remcom{yourself}.\\\hline
\end{remglyph}

\begin{prons}
\pronsrtwocone & --- & \textbf{はら (hara)}\\\hline
\rye{3}{This kun-yomi will often become voiced in second and third position when forming nouns, as in the final compound.  It never becomes voiced, however, when forming verbs.}
\pronsrthreecone & フツ (FUTSU) & --- \\\hline
\end{prons}

\begin{remprons}
\hooglig{Pay} with \comkun{hara-kiri} or \ucomon{voetsek}.\\\hline
\end{remprons}

%{\middel}
% &  \mbox{()}\\\hline
\begin{somme}
払う & \textit{pay; sweep away} & はら{\middel}う \mbox{(hara{\middel}u)}\\\hline
先払い & precede + pay = \textit{payment in advance} & さき{\middel}ばら{\middel}い \mbox{(saki{\middel}bara{\middel}i)}\\\hline
\end{somme}
\end{multicols}
\vhrulefill{5pt}
%%--------------------------------------------------------------------
%10-------------------------------------------------------------------
\begin{multicols}{2}
\begin{glyph}{Husband (夫)}{husband}\label{glyph:husband}
Husband.\\\hline
This is a good time to ensure that you can clearly distinguish the following 5 kanji: 大 ``large'' from Entry 17, 天 ``heaven'' from Entry 123, 夫 ``husband'' from here, as well as the two that included a hammer component, 失 ``lose'' from Entry 142, and 矢 ``arrow'' from Entry 285.
\end{glyph}

\begin{comps}
\comprtwocone & 一 + 大 \\\hline
\comprthreecone & One + Large \\\hline
\end{comps}

\begin{remglyph}
That \hooglig{husband} is \remcom{one} \remcom{large} \remcom{sumo wrestler}.\\\hline
\end{remglyph}

\begin{prons}
\pronsrtwocone & \textbf{フ (FU)} & \textbf{おっと (otto)}\\\hline
\pronsrthreecone & フウ (F\={U}) & --- \\\hline
\rye{3}{You are only likely to encounter this reading the final compound of Entry 504, as well as in the word 工夫 (ク{\middel}フウ/KU{\middel}F\={U}) \textit{`device, invention'}, from the kanji for `craft' and `husband' (工 being read here with its less common on-yomi).}
\end{prons}

\begin{remprons}
\hooglig{Husband} \comkun{otters} don't take \comon{foofy} \ucomon{foods}.\\\hline
\end{remprons}

%{\middel}
% &  \mbox{()}\\\hline
\begin{somme}
夫 & \textit{husband} & おっと \mbox{(otto)}\\\hline
夫人 & husband + person = \textit{wife; Mrs.} & フ{\middel}ジン \mbox{(FU{\middel}JIN)}\\\hline
\end{somme}
\end{multicols}
\vspace*{\fill}
\pagebreak
%%--------------------------------------------------------------------
%11-------------------------------------------------------------------
\vspace*{\fill}
\begin{multicols}{2}
\begin{glyph}{Alter (変)}{alter}\label{glyph:alter}
Alter, Change.\\\hline
In some compounds, the idea of an accident or mishap is present.
\end{glyph}

\begin{comps}
\comprtwocone & 夂 + 小 + 亠 \\\hline
\comprthreecone & Winter + Small + Lid/Top \\\hline
\end{comps}

\begin{remglyph}
\hooglig{Alter} under a \remcom{lid} during this \remcom{small} \remcom{winter}.\\\hline
\end{remglyph}

\begin{prons}
\pronsrtwocone & \textbf{ヘン (HEN)} & \textbf{か (ka)}\\\hline
\rye{3}{ヘン is the overwhelming choice in the first position.  The kun-yomi becomes voiced in the second position according to the same rules we learned for 付 back in Entry 226.}
\pronsrthreecone & --- & --- \\\hline
\end{prons}

\begin{remprons}
\hooglig{Alternate} between \comkun{custom} \comon{hens}.\\\hline
\end{remprons}

%{\middel}
% &  \mbox{()}\\\hline
\begin{somme}
変わる (intr.) & \textit{change; be different} & か{\middel}わる \mbox{(ka{\middel}waru)}\\\hline
変える (tr.) & \textit{alter; change} & か{\middel}える \mbox{(ka{\middel}eru)}\\\hline
変化 & alter + change = \textit{change; transformation} & ヘン{\middel}カ \mbox{(HEN{\middel}KA)}\\\hline
大変 & large + alter = \textit{serious; very} & タイ{\middel}ヘン \mbox{(TAI{\middel}HEN)}\\\hline
\end{somme}
\end{multicols}
\vhrulefill{5pt}
%%--------------------------------------------------------------------
%12-------------------------------------------------------------------
\begin{multicols}{2}
\begin{glyph}{Eat (食)}{eat}\label{glyph:eat}
Eat.
\end{glyph}

\begin{comps}
\comprtwocone & 食 \\\hline
\comprthreecone & Eat(ing) \\\hline
\end{comps}

\begin{remglyph}
Part of the Kanji radicals (\#{184}).\\\hline
\end{remglyph}

\begin{prons}
\pronsrtwocone & \textbf{ショク (SHOKU)} & \textbf{た (ta)}\\\hline
\pronsrthreecone & ジキ (JIKI) & く (ku) \\\hline
\end{prons}

\begin{remprons}
\hooglig{Eat} \comon{shockingly} in a \comkun{tunnel} in \ucomon{Tajikistan} during the \ucomkun{coup}.\\\hline
\end{remprons}

%{\middel}
% &  \mbox{()}\\\hline
\begin{somme}
食べる & \textit{eat} & た{\middel}べる \mbox{(ta{\middel}beru)}\\\hline
食べ物 & eat + thing = \textit{food} & た{\middel}べ　もの \mbox{(ta{\middel}be mono)}\\\hline
定食 & fixed + eat = \textit{set meal} & テイ{\middel}ショク \mbox{(TEI{\middel}SHOKU)}\\\hline
\end{somme}
\end{multicols}
\vspace*{\fill}
\pagebreak
%%--------------------------------------------------------------------
%13-------------------------------------------------------------------
\vspace*{\fill}
\begin{multicols}{2}
\begin{glyph}{Move (動)}{move}\label{glyph:move}
Move.
\end{glyph}

\begin{comps}
\comprtwocone & 重 + 力 \\\hline
\comprthreecone & Heavy (\ref{glyph:heavy}) + Strength/Force/Power\\\hline
\end{comps}

\begin{remglyph}
\hooglig{Move} since you are \remcom{heavy} with \remcom{strength}.\\\hline
\end{remglyph}

\begin{prons}
\pronsrtwocone & \textbf{ドウ (D\={O})} & \textbf{うご (ugo)}\\\hline
\pronsrthreecone & --- & --- \\\hline
\end{prons}

\begin{remprons}
\hooglig{Move} the \comon{door}, then can \comkun{you go}.\\\hline
\end{remprons}

%{\middel}
% &  \mbox{()}\\\hline
\begin{somme}
動く (intr.) & \textit{move} & うご{\middel}く \mbox{(ugo{\middel}ku)}\\\hline
動かす (tr.) & \textit{move} & うご{\middel}かす \mbox{(ugo{\middel}kasu)}\\\hline
動物 & move + thing = \textit{animal} & ドウ{\middel}ブツ \mbox{(D\={O}{\middel}BUTSU)}\\\hline
自動 & self + move = \textit{automatic} & ジ{\middel}ドウ \mbox{(JI{\middel}D\={O})}\\\hline
\end{somme}
\end{multicols}
\vhrulefill{5pt}
%%--------------------------------------------------------------------
%14-------------------------------------------------------------------
\begin{multicols}{2}
\begin{glyph}{Spring(s) (泉)}{springs}\label{glyph:springs}
Spring(s) (as in a source of water).
\end{glyph}

\begin{comps}
\comprtwocone & 白 + 水 \\\hline
\comprthreecone & White + Water \\\hline
\end{comps}

\begin{remglyph}
\hooglig{Springs} gushing with \remcom{white} \remcom{water}.\\\hline
\end{remglyph}

\begin{prons}
\pronsrtwocone & \textbf{セン (SEN)} & \textbf{いずみ (izumi)}\\\hline
\pronsrthreecone & --- & --- \\\hline
\end{prons}

\begin{remprons}
\hooglig{Springs} are \comon{central} to \comkun{Izumi} Kurtis' health.\\\hline
\end{remprons}

%{\middel}
% &  \mbox{()}\\\hline
\begin{somme}
泉 & \textit{spring(s)} & いずみ \mbox{(izumi)}\\\hline
\notyet{温泉} & warm + springs = \textit{hot spring; spa} & オン{\middel}セン \mbox{(ON{\middel}SEN)}\\\hline
\end{somme}
\end{multicols}
\vspace*{\fill}
\pagebreak
%%--------------------------------------------------------------------
%15-------------------------------------------------------------------
\vspace*{\fill}
\begin{multicols}{2}
\begin{glyph}{Un- (非)}{un}\label{glyph:un}
This is an important negating prefix, much like `不' (`Not', from Entry \ref{glyph:not}).  In this and future stories we will assign it the external meaning of ``amoebas''.\\\hline
Take care with the stroke order---though the first line of the `merging' component is the initial stroke of this kanji, the second comes after the three amoebas on the left.  Also keep in mind that the six amoebas are always written left to right.
\end{glyph}

\begin{comps}
\comprtwocone & 非 \\\hline
\comprthreecone & Wrong/Not \\\hline
\end{comps}

\begin{remglyph}
Part of the Kanji radicals (\#{175}).\\\hline
\end{remglyph}

\begin{prons}
\pronsrtwocone & \textbf{ヒ (HI)} & --- \\\hline
\pronsrthreecone & --- & --- \\\hline
\end{prons}

\begin{remprons}
\hooglig{Negating} \comon{Healing} powers.\\\hline
\end{remprons}

%{\middel}
% &  \mbox{()}\\\hline
\begin{somme}
非売品 & un- + sell + goods = \textit{``item not for sale''} & ヒ{\middel}バイ{\middel}ヒン \mbox{(HI{\middel}BAI{\middel}HIN)}\\\hline
\notyet{是非} &  approve + un- = \textit{without fail; by any means} & ゼ{\middel}ヒ \mbox{(ZE{\middel}HI)}\\\hline
\notyet{非公式} & un- + public + type = \textit{unofficial; informal} & ヒ{\middel}コウ{\middel}シキ \mbox{(HI{\middel}K\={O}{\middel}SHIKI)}\\\hline
\end{somme}
\end{multicols}
\vhrulefill{5pt}
%%--------------------------------------------------------------------
%16-------------------------------------------------------------------
\begin{multicols}{2}
\begin{glyph}{Substitute (代)}{substitute}\label{glyph:substitute}
Substitute.\\\hline
This interesting kanji also carries some important secondary meanings---it can indicate a `fee' for some type of good or service (think of the fee as a `substitute' for what is received), as well as convey the idea of `generations' in compounds related to eras of reigns (as in the third example below).
\end{glyph}

\begin{comps}
\comprtwocone & 弋 + 亻 \\\hline
\comprthreecone & Shoot/Ceremony + Person \\\hline
\end{comps}

\begin{remglyph}
\hooglig{Substitute} the \remcom{shooting} \remcom{person}.\\\hline
\end{remglyph}

\begin{prons}
\pronsrtwocone & \textbf{ダイ (DAI)} & \textbf{か (ka)}\\\hline
\rye{3}{Though sharing the same pronunciation, note the subtle difference in meaning between the intransitive/transitive verb pair shown here and the pair in Entry 316.}
\pronsrthreecone & タイ (TAI) & よ (yo) \\\hline
\end{prons}

\begin{remprons}
The \hooglig{substitute} \comon{died} in the \comkun{customary} way during the \ucomon{tight} \ucomkun{yomp}.\\\hline
\end{remprons}

%{\middel}
% &  \mbox{()}\\\hline
\begin{somme}
代わる (intr.) & \textit{take the place of} & か{\middel}わる \mbox{(ka{\middel}waru)}\\\hline
代える (tr.) & \textit{substitute; exchange} & か{\middel}える \mbox{(ka{\middel}eru)}\\\hline
時代 & time + substitute = \textit{era; age} & ジ{\middel}ダイ \mbox{(JI{\middel}DAI)}\\\hline
\notyet{代表}　& substitute + express = \textit{representative} & ダイ{\middel}ヒョウ \mbox{(DAI{\middel}HY\={O})}\\\hline
\end{somme}
\end{multicols}
\vspace*{\fill}
\pagebreak
%%--------------------------------------------------------------------
%17-------------------------------------------------------------------
\vspace*{\fill}
\begin{multicols}{2}
\begin{glyph}{Note (記)}{note}\label{glyph:note}
Note.\\\hline
Write down.
\end{glyph}

\begin{comps}
\comprtwocone & 言 + 己 \\\hline
\comprthreecone & Say + Oneself \\\hline
\end{comps}

\begin{remglyph}
\hooglig{Note} \remcom{oneself} while \remcom{saying} something.\\\hline
\end{remglyph}

\begin{prons}
\pronsrtwocone & \textbf{キ (KI)} & --- \\\hline
\pronsrthreecone & --- & しる (shiru) \\\hline
\end{prons}

\begin{remprons}
\hooglig{The note} said to \comon{kill} with \ucomkun{sheer oomph}.\\\hline
\end{remprons}

%{\middel}
% &  \mbox{()}\\\hline
\begin{somme}
記者 & note + individual = \textit{(newspaper) reporter / writer} & キ{\middel}シャ \mbox{(KI{\middel}SHA)}\\\hline
日記 & sun (day) + note = \textit{diary; journal} & ニッ{\middel}キ \mbox{(NIK{\middel}KI)}\\\hline
暗記 & dark + note = \textit{memorization} & アン{\middel}キ \mbox{(AN{\middel}KI)}\\\hline
\end{somme}
\end{multicols}
\vhrulefill{5pt}
%%--------------------------------------------------------------------
%18-------------------------------------------------------------------
\begin{multicols}{2}
\begin{glyph}{Reflect (映)}{reflect}\label{glyph:reflect}
Reflect an image.  Project.
\end{glyph}

\begin{comps}
\comprtwocone & 日 + 央 \\\hline
\comprthreecone & Sun + Center (\ref{glyph:center}) \\\hline
\end{comps}

\begin{remglyph}
Reflect the sun's center.\\\hline
\end{remglyph}

\begin{prons}
\pronsrtwocone & \textbf{エイ (EI)} & \textbf{うつ (utsu)}\\\hline
\pronsrthreecone & --- & は (ha) \\\hline
\end{prons}

\begin{remprons}
\hooglig{Project} your \comon{aim} at the \comkun{roots} to bring a \ucomkun{halt}.\\\hline
\end{remprons}

%{\middel}
% &  \mbox{()}\\\hline
\begin{somme}
映る (intr.) & \textit{reflect} & うつ{\middel}る \mbox{(utsu{\middel}ru)}\\\hline
移す (tr.) & \textit{reflect (something)} & うつ{\middel}す \mbox{(utsu{\middel}su)}\\\hline
映画 & reflect + painting = \textit{movie; film} & エイ{\middel}ガ \mbox{(EI{\middel}GA)}\\\hline
\end{somme}
\end{multicols}
\vspace*{\fill}
\pagebreak
%%--------------------------------------------------------------------
%19-------------------------------------------------------------------
\vspace*{\fill}
\begin{multicols}{2}
\begin{glyph}{Doctor (医)}{doctor}\label{glyph:doctor}
Doctor.  Healing.  Cure.
\end{glyph}

\begin{comps}
\comprtwocone & 匚 + 矢 \\\hline
\comprthreecone & Square Box + Arrow \\\hline
\end{comps}

\begin{remglyph}
\hooglig{Doctors} think out of the \remcom{box} at \remcom{arrow} speed.\\\hline
\end{remglyph}

\begin{prons}
\pronsrtwocone & \textbf{イ (I)} & --- \\\hline
\pronsrthreecone & --- & --- \\\hline
\end{prons}

\begin{remprons}
\hooglig{Doctors} \comon{indicate} what is wrong.\\\hline
\end{remprons}

%{\middel}
% &  \mbox{()}\\\hline
\begin{somme}
医者 & doctor + individual = \textit{doctor; physician} & イ{\middel}シャ \mbox{(I{\middel}SHA)}\\\hline
医学 & doctor + study = \textit{medical science} & イ{\middel}ガク \mbox{(I{\middel}GAKU)}\\\hline
\notyet{医院} & doctor + institution = \textit{doctor's office/clinic} & イ{\middel}イン \mbox{(I{\middel}IN)}\\\hline
\end{somme}
\end{multicols}
\vhrulefill{5pt}
%%--------------------------------------------------------------------
%20-------------------------------------------------------------------
\begin{multicols}{2}
\begin{glyph}{Ocean (洋)}{ocean}\label{glyph:ocean}
Ocean.\\\hline
This kanji is also used to denote things from the `West' (as opposed to the Orient); the first compound provides an example of this usage.
\end{glyph}

\begin{comps}
\comprtwocone & 水氵氺 + 羊 \\\hline
\comprthreecone & Water + Sheep \\\hline
\end{comps}

\begin{remglyph}
\hooglig{Ocean} is a flock of \remcom{sheep} made out of \remcom{water}.\\\hline
\end{remglyph}

\begin{prons}
\pronsrtwocone & \textbf{ヨウ (Y\={O})} & --- \\\hline
\pronsrthreecone & --- & --- \\\hline
\end{prons}

\begin{remprons}
\hooglig{Ocean} trips make you \comon{yawn} a lot.\\\hline
\end{remprons}

%{\middel}
% &  \mbox{()}\\\hline
\begin{somme}
洋服 & ocean (Western) + clothes = \textit{Western-type clothes} & ヨウ{\middel}フク \mbox{(Y\={O}{\middel}FUKU)}\\\hline
大西洋 & large + west + ocean = \textit{The Atlantic Ocean} & タイ{\middel}セイ{\middel}ヨウ \mbox{(TAI{\middel}SEI{\middel}Y\={O})}\\\hline
海洋学 & sea + ocean + study = \textit{oceanography} & カイ{\middel}ヨウ{\middel}ガク \mbox{(KAI{\middel}Y\={O}{\middel}GAKU)}\\\hline
\end{somme}
\end{multicols}
\vspace*{\fill}
\pagebreak
%%--------------------------------------------------------------------
%21-------------------------------------------------------------------
\vspace*{\fill}
\begin{multicols}{2}
\begin{glyph}{Bridge (橋)}{bridge}\label{glyph:bridge}
Bridge.
\end{glyph}

\begin{comps}
\comprtwocone & 冂 + 木 + 口 + 丿乀 + 人 \\\hline
\comprthreecone & To enclose + Tree/Wood + Mouth/Opening/Sounding + Slash + Person \\\hline
\end{comps}

\begin{remglyph}
At the \hooglig{bridge} \remcom{enclosure} a \remcom{person} \remcom{slashed} a \remcom{vampire} with \remcom{wood}.\\\hline
\end{remglyph}

\begin{prons}
\pronsrtwocone & \textbf{キョウ (KY\={O})} & \textbf{はし (hashi)}\\\hline
\rye{3}{The kun-yomi can become voiced outside of the first position (as in the final compound).}
\pronsrthreecone &  &  \\\hline
\end{prons}

\begin{remprons}
\hooglig{The bridge} in \comon{Kyoto} brought a \comkun{hushing} feeling over me.\\\hline
\end{remprons}

%{\middel}
% &  \mbox{()}\\\hline
\begin{somme}
橋 & \textit{bridge} & はし \mbox{(hashi)}\\\hline
歩道橋 & walk + road + bridge = \textit{pedestrian bridge} & ホ{\middel}ドウ{\middel}キョウ \mbox{(HO{\middel}D\={O}{\middel}KY\={O})}\\\hline
\notyet{石橋} & stone + bridge = \textit{stone bridge} & いし{\middel}ばし \mbox{(ishi{\middel}bashi)}\\\hline
\end{somme}
\end{multicols}
\vhrulefill{5pt}
%%--------------------------------------------------------------------
%22-------------------------------------------------------------------
\begin{multicols}{2}
\begin{glyph}{Present (現)}{present}\label{glyph:present}
Present (in terms of time).\\\hline
A related meaning has to do with the ideas of revealing and expressing things.\\\hline
The overall sense is of things being present or brought into being.
\end{glyph}

\begin{comps}
\comprtwocone & 王 + 見 \\\hline
\comprthreecone & King + See \\\hline
\end{comps}

\begin{remglyph}
\hooglig{Presently} the \remcom{king} \remcom{sees} everything.\\\hline
\end{remglyph}

\begin{prons}
\pronsrtwocone & \textbf{ゲン (GEN)} & \textbf{あらわ (arawa)}\\\hline
\rye{3}{The kun-yomi is the usual verb stem for the first and second examples below.  A quirk with this kanji, however, is that you may occasionally see these verbs written as 現われる and 現わす.  Think of these two verbs as things either presenting themselves (in the first instance) or being presented in some way.}
\pronsrthreecone & --- & --- \\\hline
\end{prons}

\begin{remprons}
\hooglig{Presently}, \comon{Genghis Khan} is inspecting \comkun{a raw animal}.\\\hline
\end{remprons}

%{\middel}
% &  \mbox{()}\\\hline
\begin{somme}
現れる (intr.) & \textit{appear; come out} & あらわ{\middel}れる \mbox{(arawa{\middel}reru)}\\\hline
現す (tr.) & \textit{reveal} & あらわ{\middel}す \mbox{(arawa{\middel}su)}\\\hline
現実 & present + actual = \textit{actuality; reality} & ゲン{\middel}ジツ \mbox{(GEN{\middel}JITSU)}\\\hline
\end{somme}
\end{multicols}
\vspace*{\fill}
\pagebreak
%%--------------------------------------------------------------------
%23-------------------------------------------------------------------
\vspace*{\fill}
\begin{multicols}{2}
\begin{glyph}{Ground (地)}{ground}\label{glyph:ground}
Ground.  Whereas 土 is primarily concerned with the idea of actual earth and soil, this kanji has more to do with the \textbf{surface} of things in the sense of `land', `district', and `region', for example.  This concept of a `base' is even extended to include such things as the fabric or texture of objects like cloth, as well as the background of a person's character or voice.
\end{glyph}

\begin{comps}
\comprtwocone & 土 + 也 \\\hline
\comprthreecone & Earth + To be classical/Also/Too \\\hline
\end{comps}

\begin{remglyph}
\hooglig{Grounds} were \remcom{also} \remcom{earthed}.\\\hline
\end{remglyph}

\begin{prons}
\pronsrtwocone & \textbf{チ、ジ (CHI, JI)} & --- \\\hline
\rye{3}{チ is by far the more common of these readings, especially in the second and third position, where ジ appears in only a few compounds primarily related to cloth.  In the first position the situation is a bit less precise, although even here チ will be encountered more than often.}
\pronsrthreecone & --- & --- \\\hline
\end{prons}

\begin{remprons}
\hooglig{Grounds} of a \comon{cheap} \comon{jig}.\\\hline
\end{remprons}

%{\middel}
% &  \mbox{()}\\\hline
\begin{somme}
\rye{3}{Note that the second compound uses the less common on-yomi of 土 (ト/TO); this is the only word in which the kanji has this reading.  Recall from Entry \ref{glyph:heart} that 心 in the first example below has an irregular reading of ここ/koko.}
心地 & heart + ground = \textit{feeling} & ここ{\middel}チ \mbox{(koko{\middel}CHI)}\\\hline
土地 & earth + ground = \textit{land} & ト{\middel}チ \mbox{(TO{\middel}CHI)}\\\hline
地所 & ground + location = \textit{piece of land} & ジ{\middel}ショ \mbox{(JI{\middel}SHO)}\\\hline
\end{somme}
\end{multicols}
\vhrulefill{5pt}
%%--------------------------------------------------------------------
%24-------------------------------------------------------------------
\begin{multicols}{2}
\begin{glyph}{Transmit (伝)}{transmit}\label{glyph:transmit}
Transmit.  Communicate.
\end{glyph}

\begin{comps}
\comprtwocone & 厶 + ニ + 亻 \\\hline
\comprthreecone & Self/Private + Two + Man/A Person \\\hline
\end{comps}

\begin{remglyph}
\hooglig{Transmit} \remcom{privately} between \remcom{2} \remcom{people}.\\\hline
\end{remglyph}

\begin{prons}
\pronsrtwocone & \textbf{デン (DEN)} & \textbf{つた (tsuta)}\\\hline
\rye{3}{The kun-yomi can become voiced in the second position according to the same rules for 付 (Entry \ref{glyph:attach}).}
\pronsrthreecone & テン (TEN) & --- \\\hline
\end{prons}

\begin{remprons}
\hooglig{Communicate} \comkun{stew tastes} in the \comon{den} to avoid \ucomon{tension}.\\\hline
\end{remprons}

%{\middel}
% &  \mbox{()}\\\hline
\begin{somme}
伝える & transmit; convey & つた{\middel}える \mbox{(tsuta{\middel}eru)}\\\hline
伝来 & transmit + come = \textit{transmission; handing down} & デン{\middel}ライ \mbox{(DEN{\middel}RAI)}\\\hline
伝記 & transmit + note = \textit{biography} & デン{\middel}キ \mbox{(DEN{\middel}KI)}\\\hline
\end{somme}
\end{multicols}
\vspace*{\fill}
\pagebreak
%%--------------------------------------------------------------------
%25-------------------------------------------------------------------
\vspace*{\fill}
\begin{multicols}{2}
\begin{glyph}{Bind (結)}{bind}\label{glyph:bind}
Bind, as well as the senses of `tying up' such things as a contract or marriage, for example.
\end{glyph}

\begin{comps}
\comprtwocone & 糸 + 吉 \\\hline
\comprthreecone & Thread + Lucky (\ref{glyph:lucky}) \\\hline
\end{comps}

\begin{remglyph}
\hooglig{Bind} the \remcom{lucky} \remcom{thread}.\\\hline
\end{remglyph}

\begin{prons}
\pronsrtwocone & \textbf{ケツ (KETSU)} & \textbf{むす (musu)}\\\hline
\pronsrthreecone & --- & ゆ (yu) \\\hline
\end{prons}

\begin{remprons}
\hooglig{Bound} in the \comkun{museum} is the \comon{Ketsu} \ucomkun{yule}.\\\hline
\end{remprons}

%{\middel}
% &  \mbox{()}\\\hline
\begin{somme}
結ぶ & \textit{bind; connect} & むす{\middel}ぶ \mbox{(musu{\middel}bu)}\\\hline
結婚 & bind + marriage = \textit{marriage} & ケッ{\middel}コン \mbox{(KEK{\middel}KON)}\\\hline
結合 & bind + match = \textit{union; combination} & ケツ{\middel}ゴウ \mbox{(KETSU{\middel}G\={O})}\\\hline
\end{somme}
\end{multicols}
\vhrulefill{5pt}
\vspace*{\fill}
\pagebreak
%%--------------------------------------------------------------------
